\subsection{Phân Rã Giá Trị Kỳ Dị \- Singular Value Decomposition}
\subsubsection{Giới Thiệu và Đặt Vấn Đề}

\subsubsection{Bối Cảnh và Động Lực}

Phân rã QR cho phép xử lý tốt các ma trận có hạng đầy đủ cột, nhưng nhiều bài toán
thực tế đòi hỏi phân tích sâu hơn: ma trận có thể là hình chữ nhật tùy ý, có hạng
thấp, hay bị nhiễu. \textbf{Phân rã giá trị kỳ dị} (\textit{Singular Value
Decomposition}, SVD) là công cụ mạnh mẽ nhất của đại số tuyến tính số, áp dụng
được cho \emph{mọi ma trận thực} $\mat{A} \in \R^{m \times n}$~\cite{trefethen1997}.

SVD không chỉ là một công cụ tính toán mà còn mang ý nghĩa hình học sâu sắc: nó
tiết lộ cấu trúc bên trong của một phép biến đổi tuyến tính thông qua các hướng
và độ giãn đặc trưng. Kết quả là một phân rã duy nhất (theo nghĩa chuẩn hóa) phản
ánh đầy đủ mọi tính chất đại số và hình học của ma trận.

\subsubsection{Vấn Đề Cần Giải Quyết}

\begin{tcolorbox}[enhanced, colback=blue!8!white, colframe=primaryblue,
    title=\textbf{Bài toán tổng quát}, fonttitle=\bfseries]
Cho ma trận $\mat{A} \in \R^{m \times n}$ tùy ý (không cần vuông, không cần
hạng đầy đủ). Tìm biểu diễn:
\[
    \mat{A} = \mat{U}\,\boldsymbol{\Sigma}\,\mat{V}^T,
\]
trong đó:
\begin{itemize}[noitemsep]
    \item $\mat{U} \in \R^{m \times m}$ là ma trận trực giao,
    \item $\boldsymbol{\Sigma} \in \R^{m \times n}$ là ma trận ``đường chéo'' với
          các phần tử $\sigma_1 \geq \sigma_2 \geq \cdots \geq 0$,
    \item $\mat{V} \in \R^{n \times n}$ là ma trận trực giao.
\end{itemize}
\end{tcolorbox}

\vspace{0.5em}
Từ biểu diễn này, nhiều bài toán quan trọng được giải quyết hiệu quả:
\begin{enumerate}[noitemsep]
    \item \textbf{Xấp xỉ hạng thấp:} tìm ma trận hạng $k$ gần $\mat{A}$ nhất
          (nén dữ liệu, giảm chiều).
    \item \textbf{Bình phương tối thiểu:} giải hệ $\mat{A}\mat{x} \approx \mat{b}$
          ngay cả khi $\mat{A}$ kỳ dị hoặc chữ nhật.
    \item \textbf{Phân tích thành phần chính (PCA):} tìm các hướng biến thiên lớn
          nhất trong dữ liệu.
    \item \textbf{Xử lý ảnh, hệ gợi ý, xử lý ngôn ngữ tự nhiên}, v.v.
\end{enumerate}

\subsubsection{Cơ Sở Lý Thuyết}

\subsubsection{Ma Trận Trực Giao và Tính Chất}

\begin{mydef}{Ma trận trực giao}{}
Ma trận $\mat{Q} \in \R^{n \times n}$ là \textbf{trực giao} nếu
$\mat{Q}^T\mat{Q} = \mat{Q}\mat{Q}^T = \mat{I}_n$, tức $\mat{Q}^{-1} = \mat{Q}^T$.
\end{mydef}

\begin{mythm}{Bảo toàn chuẩn}{}
Với ma trận trực giao $\mat{Q}$: $\norm{\mat{Q}\mat{x}}_2 = \norm{\mat{x}}_2$ với
mọi $\mat{x}$.
\end{mythm}

\begin{proof}
$\norm{\mat{Q}\mat{x}}_2^2 = (\mat{Q}\mat{x})^T(\mat{Q}\mat{x}) =
\mat{x}^T\underbrace{\mat{Q}^T\mat{Q}}_{=\,\mat{I}}\mat{x} = \norm{\mat{x}}_2^2$.
\end{proof}

\subsubsection{Ma Trận Đối Xứng Xác Định Không Âm}

\begin{mydef}{Ma trận đối xứng xác định không âm (SPSD)}{}
Ma trận $\mat{S} \in \R^{n \times n}$ là \textbf{đối xứng xác định không âm}
(\textit{Symmetric Positive SemiDefinite}, SPSD) nếu $\mat{S}^T = \mat{S}$ và
$\mat{x}^T\mat{S}\mat{x} \geq 0$ với mọi $\mat{x} \in \R^n$.
\end{mydef}

\begin{myprop}{$\mat{A}^T\mat{A}$ là SPSD}{}
Với mọi $\mat{A} \in \R^{m \times n}$, ma trận $\mat{A}^T\mat{A}$ là đối xứng
xác định không âm.
\end{myprop}

\begin{proof}
\textit{Đối xứng:} $(\mat{A}^T\mat{A})^T = \mat{A}^T(\mat{A}^T)^T = \mat{A}^T\mat{A}$. \\
\textit{Xác định không âm:} $\mat{x}^T(\mat{A}^T\mat{A})\mat{x} =
(\mat{A}\mat{x})^T(\mat{A}\mat{x}) = \norm{\mat{A}\mat{x}}_2^2 \geq 0$.
\end{proof}

\begin{mythm}{Định lý phổ cho ma trận đối xứng}{}
Mọi ma trận đối xứng $\mat{S} \in \R^{n \times n}$ đều có thể chéo hóa trực giao:
$\mat{S} = \mat{V}\boldsymbol{\Lambda}\mat{V}^T$ với $\mat{V}$ trực giao và
$\boldsymbol{\Lambda} = \operatorname{diag}(\lambda_1, \ldots, \lambda_n)$ chứa
các giá trị riêng thực. Nếu $\mat{S}$ là SPSD thì $\lambda_i \geq 0$ với mọi $i$.
\end{mythm}

\subsubsection{Định Nghĩa Giá Trị Kỳ Dị}

\begin{mydef}{Giá trị kỳ dị và vectơ kỳ dị}{}
Cho $\mat{A} \in \R^{m \times n}$. Áp dụng định lý phổ cho $\mat{A}^T\mat{A}$ với
các giá trị riêng $\lambda_1 \geq \lambda_2 \geq \cdots \geq \lambda_n \geq 0$:
\begin{itemize}[noitemsep]
    \item \textbf{Giá trị kỳ dị:} $\sigma_i = \sqrt{\lambda_i(\mat{A}^T\mat{A})}$,
          thường sắp xếp $\sigma_1 \geq \sigma_2 \geq \cdots \geq \sigma_n \geq 0$.
    \item \textbf{Vectơ kỳ dị phải:} các cột $\mat{v}_i$ của ma trận $\mat{V}$
          trong phân tích $\mat{A}^T\mat{A} = \mat{V}\boldsymbol{\Lambda}\mat{V}^T$.
    \item \textbf{Vectơ kỳ dị trái:} $\mat{u}_i = \frac{1}{\sigma_i}\mat{A}\mat{v}_i$
          (với $\sigma_i > 0$).
\end{itemize}
\end{mydef}

\subsubsection{Phân Rã Giá Trị Kỳ Dị — Định Lý và Chứng Minh}

\subsubsection{Phát Biểu Định Lý}

\begin{mythm}{Phân rã SVD (Full SVD)}{svd_exist}
Với mọi $\mat{A} \in \R^{m \times n}$, tồn tại các ma trận trực giao
$\mat{U} \in \R^{m \times m}$, $\mat{V} \in \R^{n \times n}$ và ma trận
$\boldsymbol{\Sigma} \in \R^{m \times n}$ với $\Sigma_{ij} = 0$ khi $i \neq j$
và $\sigma_1 \geq \sigma_2 \geq \cdots \geq \sigma_p \geq 0$ ($p = \min(m,n)$)
trên đường chéo, sao cho:
\[
    \mat{A} = \mat{U}\,\boldsymbol{\Sigma}\,\mat{V}^T.
\]
Các $\sigma_i$ là duy nhất (gọi là \textbf{giá trị kỳ dị} của $\mat{A}$).
\end{mythm}

\begin{mydef}{SVD gọn (Thin/Economy SVD)}{}
Nếu $r = \rank(\mat{A})$, định nghĩa:
$\mat{U}_r = [\mat{u}_1 \mid \cdots \mid \mat{u}_r] \in \R^{m \times r}$,
$\boldsymbol{\Sigma}_r = \operatorname{diag}(\sigma_1, \ldots, \sigma_r) \in \R^{r \times r}$,
$\mat{V}_r = [\mat{v}_1 \mid \cdots \mid \mat{v}_r] \in \R^{n \times r}$.
Khi đó \textbf{SVD gọn} của $\mat{A}$ là:
\[
    \mat{A} = \mat{U}_r\,\boldsymbol{\Sigma}_r\,\mat{V}_r^T.
\]
Dạng này tiết kiệm hơn vì chỉ lưu $r$ cột/hàng thay vì toàn bộ $m \times n$.
\end{mydef}

\begin{note}[Full SVD so với Thin SVD]
Full SVD giữ toàn bộ ma trận trực giao $\mat{U}$ và $\mat{V}$, bao gồm cả các
cột ứng với $\sigma_i = 0$. Thin SVD chỉ giữ $r$ cột hữu ích.
Trong tính toán, Thin SVD thường được dùng vì hiệu quả hơn.
\end{note}

\subsubsection{Chứng Minh Định Lý}

\begin{proof}
Chứng minh bằng cách \emph{xây dựng tường minh} $\mat{U}$, $\boldsymbol{\Sigma}$,
$\mat{V}$~\cite{golub2013}.

\medskip
\noindent\textbf{Bước 1: Phân tích $\mat{A}^T\mat{A}$ và định nghĩa các $\sigma_i$.}

Theo Mệnh đề~2.1, $\mat{A}^T\mat{A}$ là SPSD. Theo định lý phổ:
\[
    \mat{A}^T\mat{A} = \mat{V}\boldsymbol{\Lambda}\mat{V}^T,
\]
với $\mat{V} \in \R^{n \times n}$ trực giao và $\boldsymbol{\Lambda} =
\operatorname{diag}(\lambda_1, \ldots, \lambda_n)$, $\lambda_1 \geq \cdots \geq \lambda_r
> 0 = \lambda_{r+1} = \cdots = \lambda_n$ ($r = \rank(\mat{A})$).
Đặt $\sigma_i = \sqrt{\lambda_i}$ cho $i = 1, \ldots, r$.

\medskip
\noindent\textbf{Bước 2: Xây dựng các vectơ kỳ dị trái $\mat{u}_i$ ($i \leq r$).}

Với mỗi $i = 1, \ldots, r$, đặt:
\[
    \mat{u}_i \;=\; \frac{1}{\sigma_i}\mat{A}\mat{v}_i.
\]

\medskip
\noindent\textbf{Bước 3: $\{\mat{u}_1, \ldots, \mat{u}_r\}$ là tập trực chuẩn.}

Với $1 \leq i, j \leq r$:
\[
    \inner{\mat{u}_i}{\mat{u}_j}
    = \frac{1}{\sigma_i\sigma_j}\mat{v}_i^T\mat{A}^T\mat{A}\mat{v}_j
    = \frac{1}{\sigma_i\sigma_j}\mat{v}_i^T\underbrace{(\mat{A}^T\mat{A}\mat{v}_j)}_{=\,\lambda_j\mat{v}_j}
    = \frac{\lambda_j}{\sigma_i\sigma_j}\underbrace{\mat{v}_i^T\mat{v}_j}_{=\,\delta_{ij}}.
\]
Khi $i = j$: $\lambda_i/\sigma_i^2 = 1$. \quad Khi $i \neq j$: $0$.
Vậy $\inner{\mat{u}_i}{\mat{u}_j} = \delta_{ij}$.

\medskip
\noindent\textbf{Bước 4: Với $i > r$, $\mat{A}\mat{v}_i = \mat{0}$.}

Ta cần bổ đề: $\operatorname{null}(\mat{A}^T\mat{A}) = \operatorname{null}(\mat{A})$.
\begin{itemize}[noitemsep, topsep=2pt]
    \item[$(\subseteq)$] Nếu $\mat{A}\mat{x} = \mat{0}$ thì $\mat{A}^T\mat{A}\mat{x} = \mat{0}$.
    \item[$(\supseteq)$] Nếu $\mat{A}^T\mat{A}\mat{x} = \mat{0}$ thì $\norm{\mat{A}\mat{x}}_2^2 = \mat{x}^T\mat{A}^T\mat{A}\mat{x} = 0$, suy ra $\mat{A}\mat{x} = \mat{0}$.
\end{itemize}
Vì $i > r$, $\lambda_i = 0$ nên $\mat{v}_i \in \operatorname{null}(\mat{A}^T\mat{A}) =
\operatorname{null}(\mat{A})$, tức $\mat{A}\mat{v}_i = \mat{0}$.

\medskip
\noindent\textbf{Bước 5: Mở rộng thành cơ sở trực chuẩn của $\R^m$.}

$\{\mat{u}_1, \ldots, \mat{u}_r\}$ là tập trực chuẩn trong $\R^m$. Mở rộng bằng
Gram--Schmidt thành cơ sở trực chuẩn đầy đủ $\{\mat{u}_1, \ldots, \mat{u}_m\}$.
Đặt $\mat{U} = [\mat{u}_1 \mid \cdots \mid \mat{u}_m] \in \R^{m \times m}$ (trực giao).

\medskip
\noindent\textbf{Bước 6: Xây dựng $\boldsymbol{\Sigma}$ và kiểm tra $\mat{A} = \mat{U}\boldsymbol{\Sigma}\mat{V}^T$.}

Đặt $\boldsymbol{\Sigma} \in \R^{m \times n}$ với $\Sigma_{ii} = \sigma_i$ ($i = 1,\ldots,r$),
các phần tử còn lại bằng 0. Xét tích $\mat{A}\mat{V}$:
\[
    \mat{A}\mat{V} = [\mat{A}\mat{v}_1 \mid \cdots \mid \mat{A}\mat{v}_n]
    = [\sigma_1\mat{u}_1 \mid \cdots \mid \sigma_r\mat{u}_r \mid \mat{0} \mid \cdots \mid \mat{0}]
    = \mat{U}\boldsymbol{\Sigma}.
\]
Vì $\mat{V}$ trực giao ($\mat{V}^{-1} = \mat{V}^T$):
\[
    \mat{A} = \mat{U}\boldsymbol{\Sigma}\mat{V}^T.
\]

\textbf{Tính duy nhất của $\sigma_i$:} Vì $\mat{A}^T\mat{A} = \mat{V}\boldsymbol{\Sigma}^T
\boldsymbol{\Sigma}\mat{V}^T$ và $\boldsymbol{\Sigma}^T\boldsymbol{\Sigma} =
\operatorname{diag}(\sigma_1^2,\ldots,\sigma_r^2,0,\ldots,0)$, các $\sigma_i^2$ là
các giá trị riêng dương của $\mat{A}^T\mat{A}$, xác định duy nhất. Do đó $\sigma_i$
duy nhất.
\end{proof}

\begin{mycor}{Liên hệ với eigendecomposition}{}
Từ $\mat{A} = \mat{U}\boldsymbol{\Sigma}\mat{V}^T$, ta suy ra hai phân tích chéo:
\begin{align*}
    \mat{A}^T\mat{A} &= \mat{V}\,(\boldsymbol{\Sigma}^T\boldsymbol{\Sigma})\,\mat{V}^T
    \quad \text{(giá trị riêng của $\mat{A}^T\mat{A}$ là $\sigma_i^2$)},\\
    \mat{A}\mat{A}^T &= \mat{U}\,(\boldsymbol{\Sigma}\boldsymbol{\Sigma}^T)\,\mat{U}^T
    \quad \text{(giá trị riêng khác 0 của $\mat{A}\mat{A}^T$ cũng là $\sigma_i^2$)}.
\end{align*}
Đây là cơ sở của các thuật toán tính SVD qua eigendecomposition~\cite{trefethen1997}.
\end{mycor}

\subsubsection{Tính Chất Hình Học của SVD}

\subsubsection{SVD = Quay -- Co Giãn -- Quay}

Mọi phép biến đổi tuyến tính $\mat{x} \mapsto \mat{A}\mat{x}$ đều có thể phân
tích thành ba bước kế tiếp~\cite{strang2023}:
\[
    \mat{x}
    \;\xrightarrow{\quad\mat{V}^T\quad}\;
    \mat{V}^T\mat{x}
    \;\xrightarrow{\quad\boldsymbol{\Sigma}\quad}\;
    \boldsymbol{\Sigma}\mat{V}^T\mat{x}
    \;\xrightarrow{\quad\mat{U}\quad}\;
    \mat{U}\boldsymbol{\Sigma}\mat{V}^T\mat{x} = \mat{A}\mat{x}.
\]
\begin{itemize}[noitemsep]
    \item $\mat{V}^T$: \textbf{phép quay/phản xạ} trong $\R^n$ (chuyển cơ sở
          theo $\{\mat{v}_i\}$).
    \item $\boldsymbol{\Sigma}$: \textbf{co giãn} theo các trục tọa độ với tỉ lệ
          $\sigma_1, \sigma_2, \ldots$ (và có thể thay đổi chiều $n \to m$).
    \item $\mat{U}$: \textbf{phép quay/phản xạ} trong $\R^m$ (chuyển sang cơ sở
          $\{\mat{u}_i\}$).
\end{itemize}

\subsubsection{Hình Ảnh của Hình Cầu Đơn Vị}

\begin{mythm}{Hình cầu đơn vị $\to$ Hyperellipsoid}{}
Tập $\{\mat{A}\mat{x} : \norm{\mat{x}}_2 = 1\}$ (hình ảnh của hình cầu đơn vị
qua $\mat{A}$) là một hyperellipsoid với các bán trục $\sigma_1 \geq \sigma_2 \geq
\cdots \geq \sigma_r > 0$ theo các hướng $\mat{u}_1, \mat{u}_2, \ldots, \mat{u}_r$.
\end{mythm}

\begin{proof}
Đặt $\mat{y} = \mat{V}^T\mat{x}$. Vì $\mat{V}$ trực giao, $\norm{\mat{x}}_2 = 1
\Leftrightarrow \norm{\mat{y}}_2 = 1$.
Tập $\{\mat{U}\boldsymbol{\Sigma}\mat{y} : \norm{\mat{y}}_2 = 1\}$ bằng tập
$\{\mat{U}\mat{z} : \mat{z} = \boldsymbol{\Sigma}\mat{y},\, \norm{\mat{y}}_2 = 1\}$.
Với $\mat{y} = (y_1, \ldots, y_n)^T$ và $\norm{\mat{y}}_2 = 1$:
\[
    \mat{z} = \boldsymbol{\Sigma}\mat{y}
    \implies z_i = \sigma_i y_i,
    \implies \sum_{i=1}^r \frac{z_i^2}{\sigma_i^2} = \sum_{i=1}^r y_i^2 \leq 1.
\]
Tập $\{\mat{z}\}$ là ellipsoid căn chuẩn. Nhân với $\mat{U}$ (phép quay) cho
ellipsoid với bán trục $\sigma_i$ theo hướng $\mat{u}_i$.
\end{proof}

\begin{figure}[H]
\centering
\begin{tikzpicture}[scale=0.92, >=Stealth]
  % Panel 1: Input space (unit circle)
  \begin{scope}[xshift=0cm]
    \draw[primaryblue!15, fill=primaryblue!6, thick, primaryblue]
      (0,0) circle(1);
    % Grid lines
    \draw[gray!35, thin] (-1.4,0)--(1.4,0);
    \draw[gray!35, thin] (0,-1.4)--(0,1.4);
    % v1 at 45 degrees
    \draw[->, very thick, primaryblue] (0,0)
      -- ({cos(45)},{sin(45)})
      node[above right, primaryblue, font=\small] {$\mathbf{v}_1$};
    % v2 at 135 degrees
    \draw[->, very thick, accentblue] (0,0)
      -- ({cos(135)},{sin(135)})
      node[above left, accentblue, font=\small] {$\mathbf{v}_2$};
    \fill[black] (0,0) circle(1.5pt);
    \node[font=\small, below] at (0,-1.65)
      {\textit{Đầu vào}: $\|\mathbf{x}\|=1$};
  \end{scope}

  % Arrow label (V^T then Sigma)
  \draw[->, very thick, darkgray] (1.2,0) -- (2.1,0);
  \node[above, darkgray, font=\footnotesize] at (1.65, 0.12)
    {$\boldsymbol{\Sigma}\mathbf{V}^T$};

  % Panel 2: After Sigma*V^T (axis-aligned ellipse)
  \begin{scope}[xshift=3.8cm]
    \draw[successgreen!15, fill=successgreen!6, thick, successgreen!70!black]
      (0,0) ellipse(1.6 and 0.55);
    % Grid lines
    \draw[gray!35, thin] (-1.9,0)--(1.9,0);
    \draw[gray!35, thin] (0,-0.85)--(0,0.85);
    % sigma_1 annotation
    \draw[<->, successgreen!70!black, thick] (0,-0.78) -- (1.6,-0.78);
    \node[below, successgreen!70!black, font=\footnotesize] at (0.8,-0.78)
      {$\sigma_1$};
    % sigma_2 annotation
    \draw[<->, successgreen!70!black, thick] (1.85,0) -- (1.85,0.55);
    \node[right, successgreen!70!black, font=\footnotesize] at (1.9,0.27)
      {$\sigma_2$};
    \fill[black] (0,0) circle(1.5pt);
    \node[font=\small, below] at (0,-1.65)
      {\textit{Sau} $\boldsymbol{\Sigma}\mathbf{V}^T$};
  \end{scope}

  % Arrow label U
  \draw[->, very thick, darkgray] (5.65,0) -- (6.55,0);
  \node[above, darkgray, font=\footnotesize] at (6.1, 0.12)
    {$\mathbf{U}$};

  % Panel 3: Output space (rotated ellipse)
  \begin{scope}[xshift=8.5cm]
    \draw[warnorange!15, fill=warnorange!6, thick, warnorange, rotate=30]
      (0,0) ellipse(1.6 and 0.55);
    % Grid lines
    \draw[gray!35, thin] (-1.9,0)--(1.9,0);
    \draw[gray!35, thin] (0,-1.0)--(0,1.0);
    % u1 at 30 deg, length = sigma_1 = 1.6
    \draw[->, very thick, primaryblue] (0,0)
      -- ({1.6*cos(30)},{1.6*sin(30)})
      node[above right, primaryblue, font=\small] {$\sigma_1\mathbf{u}_1$};
    % u2 at 120 deg, length = sigma_2 = 0.55
    \draw[->, very thick, accentblue] (0,0)
      -- ({0.55*cos(120)},{0.55*sin(120)})
      node[above left, accentblue, font=\small] {$\sigma_2\mathbf{u}_2$};
    \fill[black] (0,0) circle(1.5pt);
    \node[font=\small, below] at (0,-1.65)
      {\textit{Đầu ra}: $\mathbf{Ax}$};
  \end{scope}
\end{tikzpicture}
\caption{Tính chất hình học của SVD trong $\R^2$: $\mat{A} = \mat{U}\boldsymbol{\Sigma}\mat{V}^T$
biến đổi hình tròn đơn vị (trái, \textcolor{primaryblue}{xanh}) thành ellipse
theo trục tọa độ (giữa, \textcolor{successgreen!70!black}{xanh lá}) bằng
$\boldsymbol{\Sigma}\mat{V}^T$, rồi quay thành ellipse cuối (phải,
\textcolor{warnorange}{cam}) bằng $\mat{U}$. Bán trục dài/ngắn là $\sigma_1$,
$\sigma_2$ theo các hướng $\mat{u}_1$, $\mat{u}_2$.}
\label{fig:svd_geometry}
\end{figure}

\subsubsection{Tính Chất Đại Số}

\subsubsection{Rank và Các Không Gian Con}

\begin{mythm}{Rank qua SVD}{}
Cho $\mat{A} = \mat{U}\boldsymbol{\Sigma}\mat{V}^T$ với $r$ giá trị kỳ dị dương.
Khi đó $\rank(\mat{A}) = r$ và:
\begin{align*}
    \col(\mat{A}) &= \spann\{\mat{u}_1, \ldots, \mat{u}_r\}, \\
    \operatorname{null}(\mat{A}) &= \spann\{\mat{v}_{r+1}, \ldots, \mat{v}_n\}, \\
    \operatorname{row}(\mat{A}) &= \spann\{\mat{v}_1, \ldots, \mat{v}_r\}, \\
    \operatorname{null}(\mat{A}^T) &= \spann\{\mat{u}_{r+1}, \ldots, \mat{u}_m\}.
\end{align*}
\end{mythm}

\begin{proof}
Từ $\mat{A}\mat{v}_i = \sigma_i\mat{u}_i$ (với $\sigma_i > 0$ khi $i \leq r$ và
$= 0$ khi $i > r$): các cột của $\mat{A}$ nằm trong $\spann\{\mat{u}_1,\ldots,\mat{u}_r\}$,
và mỗi $\mat{u}_i$ ($i \leq r$) được biểu diễn qua cột của $\mat{A}$, nên
$\col(\mat{A}) = \spann\{\mat{u}_1,\ldots,\mat{u}_r\}$, suy ra $\rank(\mat{A}) = r$.
Tương tự cho các không gian còn lại.
\end{proof}

\subsubsection{Chuẩn Ma Trận}

\begin{mythm}{Chuẩn spectral và Frobenius qua SVD}{svd_norms}
Cho $\mat{A} = \mat{U}\boldsymbol{\Sigma}\mat{V}^T$ với $r$ giá trị kỳ dị dương:
\[
    \norm{\mat{A}}_2 = \sigma_1,
    \qquad
    \norm{\mat{A}}_F = \sqrt{\sigma_1^2 + \sigma_2^2 + \cdots + \sigma_r^2}.
\]
\end{mythm}

\begin{proof}
\textit{Chuẩn spectral:}
$\norm{\mat{A}}_2 = \max_{\norm{\mat{x}}_2=1}\norm{\mat{A}\mat{x}}_2 =
\max_{\norm{\mat{y}}_2=1}\norm{\boldsymbol{\Sigma}\mat{y}}_2$
(đặt $\mat{y} = \mat{V}^T\mat{x}$, dùng tính bảo toàn chuẩn của $\mat{U}$ và $\mat{V}$).
$\norm{\boldsymbol{\Sigma}\mat{y}}_2^2 = \sum_{i=1}^r \sigma_i^2 y_i^2 \leq \sigma_1^2\sum y_i^2
= \sigma_1^2$, đẳng thức khi $\mat{y} = \mat{e}_1$. Suy ra $\norm{\mat{A}}_2 = \sigma_1$.

\textit{Chuẩn Frobenius:}
$\norm{\mat{A}}_F^2 = \operatorname{tr}(\mat{A}^T\mat{A}) =
\operatorname{tr}(\mat{V}\boldsymbol{\Sigma}^T\boldsymbol{\Sigma}\mat{V}^T) =
\operatorname{tr}(\boldsymbol{\Sigma}^T\boldsymbol{\Sigma}) = \sum_{i=1}^r \sigma_i^2$.
\end{proof}

\subsubsection{Dạng Tích Ngoài (Outer Product Form)}

\begin{mythm}{Biểu diễn tổng tích ngoài}{}
\[
    \mat{A} = \sum_{i=1}^{r} \sigma_i\, \mat{u}_i\mat{v}_i^T.
\]
Đây là tổng $r$ ma trận hạng 1, mỗi ma trận đóng góp một ``lớp thông tin'' với trọng
số $\sigma_i$.
\end{mythm}

\begin{proof}
Từ $\mat{A} = \mat{U}_r\boldsymbol{\Sigma}_r\mat{V}_r^T$, khai triển theo định nghĩa
tích ma trận theo cột:
$\mat{U}_r\boldsymbol{\Sigma}_r = [\sigma_1\mat{u}_1 \mid \cdots \mid \sigma_r\mat{u}_r]$,
và $(\mat{U}_r\boldsymbol{\Sigma}_r)\mat{V}_r^T = \sum_{i=1}^r \sigma_i\mat{u}_i\mat{v}_i^T$.
\end{proof}

\subsubsection{Giả Nghịch Đảo Moore--Penrose}

\begin{mydef}{Giả nghịch đảo Moore--Penrose}{}
Giả nghịch đảo (\textit{pseudoinverse}) của $\mat{A} = \mat{U}\boldsymbol{\Sigma}\mat{V}^T$
là:
\[
    \mat{A}^+ = \mat{V}\,\boldsymbol{\Sigma}^+\,\mat{U}^T,
\]
trong đó $\boldsymbol{\Sigma}^+ \in \R^{n \times m}$ có
$\Sigma^+_{ii} = 1/\sigma_i$ với $\sigma_i > 0$, và 0 ở các vị trí khác.
Dạng tường minh: $\mat{A}^+ = \mat{V}_r\boldsymbol{\Sigma}_r^{-1}\mat{U}_r^T$.
\end{mydef}

$\mat{A}^+$ thỏa bốn điều kiện Moore--Penrose: $\mat{A}\mat{A}^+\mat{A} = \mat{A}$,
$\mat{A}^+\mat{A}\mat{A}^+ = \mat{A}^+$, $(\mat{A}\mat{A}^+)^T = \mat{A}\mat{A}^+$,
$(\mat{A}^+\mat{A})^T = \mat{A}^+\mat{A}$. Khi $\mat{A}$ vuông khả nghịch:
$\mat{A}^+ = \mat{A}^{-1}$.

\subsubsection{Ứng Dụng}

\subsubsection{Xấp Xỉ Hạng Thấp --- Định Lý Eckart--Young}

\subsubsubsection{Định Nghĩa Xấp Xỉ Hạng $k$}

\begin{mydef}{Xấp xỉ hạng $k$ tốt nhất}{}
Với $1 \leq k \leq r = \rank(\mat{A})$, định nghĩa:
\[
    \mat{A}_k \;=\; \sum_{i=1}^{k} \sigma_i\,\mat{u}_i\mat{v}_i^T
    \;=\; \mat{U}_k\,\boldsymbol{\Sigma}_k\,\mat{V}_k^T,
\]
là ma trận hạng $k$ thu được bằng cách chỉ giữ lại $k$ số hạng lớn nhất trong
dạng tích ngoài của $\mat{A}$.
\end{mydef}

\begin{figure}[H]
\centering
\begin{tikzpicture}[scale=1.0, >=Stealth]
  % Axes
  \draw[->, darkgray, thick] (-0.4, 0) -- (8.2, 0) node[right] {$i$};
  \draw[->, darkgray, thick] (0, -0.2) -- (0, 4.0) node[above] {$\sigma_i$};

  % Bars (decreasing singular values)
  \foreach \i/\h/\clr in {
      1/3.60/primaryblue,
      2/2.90/primaryblue,
      3/1.80/primaryblue,
      4/0.75/black!45,
      5/0.40/black!35,
      6/0.20/black!28,
      7/0.08/black!22}
  {
    \fill[\clr] (\i*1.0-0.35, 0) rectangle (\i*1.0+0.35, \h);
    \draw[black!40] (\i*1.0-0.35, 0) rectangle (\i*1.0+0.35, \h);
    \node[below, font=\footnotesize, black!60] at (\i*1.0, -0.05) {$\i$};
  }

  % Rank-k cutoff at k=3 (between bar 3 and bar 4)
  \draw[warnorange, dashed, very thick] (3.5, -0.3) -- (3.5, 4.0);
  \node[warnorange, font=\small, rotate=90, anchor=south] at (3.65, 1.5)
    {Giữ lại $k=3$};

  % Energy annotation
  \draw[<->, successgreen!70!black, thick] (0.05, 3.75) -- (3.48, 3.75);
  \node[successgreen!70!black, font=\footnotesize, above] at (1.75, 3.75)
    {năng lượng giữ lại};
  \draw[<->, accentblue, thick] (3.52, 3.75) -- (7.35, 3.75);
  \node[accentblue, font=\footnotesize, above] at (5.4, 3.75)
    {năng lượng bị loại};

  % Sigma_k+1 label
  \draw[warnorange, <-, thick] (4.0, 0.75) -- (5.2, 1.3);
  \node[warnorange, font=\footnotesize, anchor=west] at (5.2, 1.3)
    {$\sigma_{k+1}$ = sai số tối ưu};

  % Axis labels
  \node[darkgray, font=\small] at (4.0, -0.65) {Chỉ số $i$};
\end{tikzpicture}
\caption{Phổ giá trị kỳ dị của $\mat{A}$. Các thanh \textcolor{primaryblue}{xanh}
($i \leq k=3$) là các thành phần được giữ lại trong xấp xỉ $\mat{A}_k$;
các thanh \textcolor{darkgray!60}{xám} bị loại bỏ. Sai số xấp xỉ tốt nhất
(chuẩn spectral) đúng bằng $\sigma_{k+1}$.}
\label{fig:svd_spectrum}
\end{figure}

\subsubsubsection{Định Lý Eckart--Young--Mirsky}

\begin{mythm}{Eckart--Young--Mirsky}{}
Cho $\mat{A} \in \R^{m \times n}$ với $\rank(\mat{A}) = r > k$. Với chuẩn spectral
và chuẩn Frobenius:
\[
    \min_{\substack{\mat{B} \in \R^{m \times n} \\ \rank(\mat{B}) \leq k}}
    \norm{\mat{A} - \mat{B}}_2 = \norm{\mat{A} - \mat{A}_k}_2 = \sigma_{k+1},
\]
\[
    \min_{\substack{\mat{B} \in \R^{m \times n} \\ \rank(\mat{B}) \leq k}}
    \norm{\mat{A} - \mat{B}}_F = \norm{\mat{A} - \mat{A}_k}_F
    = \sqrt{\sigma_{k+1}^2 + \cdots + \sigma_r^2}.
\]
Nghiệm tối ưu là $\mat{A}_k$ trong cả hai trường hợp~\cite{golub2013}.
\end{mythm}

\begin{proof}
Chứng minh cho chuẩn spectral (trường hợp Frobenius tương tự).

\medskip
\noindent\textit{Phần 1: Tính sai số của $\mat{A} - \mat{A}_k$.}

$\mat{A} - \mat{A}_k = \displaystyle\sum_{i=k+1}^{r}\sigma_i\mat{u}_i\mat{v}_i^T$
là SVD của ma trận phần còn lại. Theo Định lý~\ref{thm:svd_norms}:
$\norm{\mat{A} - \mat{A}_k}_2 = \sigma_{k+1}$.

\medskip
\noindent\textit{Phần 2: Mọi ma trận hạng $k$ đều có sai số $\geq \sigma_{k+1}$.}

Giả sử $\mat{B}$ có $\rank(\mat{B}) \leq k$. Khi đó:
\[
    \dim\bigl(\operatorname{null}(\mat{B})\bigr) \geq n - k.
\]
Xét không gian con $\mathcal{W} = \spann\{\mat{v}_1, \ldots, \mat{v}_{k+1}\}$
với $\dim(\mathcal{W}) = k+1$.

Theo bất đẳng thức chiều:
\[
    \dim\bigl(\operatorname{null}(\mat{B}) \cap \mathcal{W}\bigr)
    \geq \underbrace{(n-k)}_{\dim\operatorname{null}(\mat{B})}
    + \underbrace{(k+1)}_{\dim\mathcal{W}} - n = 1.
\]
Suy ra tồn tại vectơ đơn vị $\mat{w} \in \operatorname{null}(\mat{B}) \cap \mathcal{W}$.
Viết $\mat{w} = \sum_{i=1}^{k+1}\alpha_i\mat{v}_i$ với $\sum\alpha_i^2 = 1$.
Vì $\mat{B}\mat{w} = \mat{0}$:
\begin{align*}
    \norm{\mat{A} - \mat{B}}_2
    &\geq \norm{(\mat{A} - \mat{B})\mat{w}}_2
    = \norm{\mat{A}\mat{w}}_2 \\
    &= \left\|\sum_{i=1}^{k+1}\alpha_i\mat{A}\mat{v}_i\right\|_2
    = \left\|\sum_{i=1}^{k+1}\alpha_i\sigma_i\mat{u}_i\right\|_2 \\
    &= \sqrt{\sum_{i=1}^{k+1}\alpha_i^2\sigma_i^2}
    \;\geq\; \sigma_{k+1}\sqrt{\sum_{i=1}^{k+1}\alpha_i^2}
    = \sigma_{k+1}.
\end{align*}
Kết hợp hai phần: $\displaystyle\min_{\rank(\mat{B})\leq k}\norm{\mat{A}-\mat{B}}_2 = \sigma_{k+1}$,
đạt được khi $\mat{B} = \mat{A}_k$.
\end{proof}

\begin{note}[Ý nghĩa thực tế]
Định lý Eckart--Young cho biết: \textbf{không có ma trận hạng $k$ nào xấp xỉ $\mat{A}$
tốt hơn $\mat{A}_k$} (theo cả chuẩn spectral lẫn Frobenius). Đây là cơ sở lý thuyết
của nén ảnh, PCA, LSA (\textit{Latent Semantic Analysis}), và các phương pháp
giảm chiều dữ liệu.
\end{note}


\subsubsection{Giải Hệ Bình Phương Tối Thiểu qua SVD}

Xét bài toán $\min_{\mat{x}}\norm{\mat{A}\mat{x} - \mat{b}}_2$ với
$\mat{A} \in \R^{m \times n}$, $\mat{b} \in \R^m$ tùy ý (không yêu cầu
$\mat{A}$ có hạng đầy đủ cột).

\begin{mythm}{Nghiệm qua giả nghịch đảo}{}
Nghiệm bình phương tối thiểu có chuẩn nhỏ nhất là~\cite{strang2023}:
\[
    \hat{\mat{x}} = \mat{A}^+\mat{b}
    = \mat{V}_r\boldsymbol{\Sigma}_r^{-1}\mat{U}_r^T\mat{b}
    = \sum_{i=1}^{r} \frac{\inner{\mat{u}_i}{\mat{b}}}{\sigma_i}\,\mat{v}_i.
\]
\end{mythm}

\begin{proof}[Ý tưởng chứng minh]
Phân tích $\mat{b} = \mat{U}\mat{U}^T\mat{b}$ và dùng tính bảo toàn chuẩn:
\begin{align*}
    \norm{\mat{A}\mat{x} - \mat{b}}_2^2
    &= \norm{\mat{U}\boldsymbol{\Sigma}\mat{V}^T\mat{x} - \mat{b}}_2^2
    = \norm{\boldsymbol{\Sigma}\mat{V}^T\mat{x} - \mat{U}^T\mat{b}}_2^2.
\end{align*}
Đặt $\mat{y} = \mat{V}^T\mat{x}$ và $\mat{c} = \mat{U}^T\mat{b}$. Khai triển:
\[
    \norm{\boldsymbol{\Sigma}\mat{y} - \mat{c}}_2^2
    = \sum_{i=1}^r(\sigma_i y_i - c_i)^2 + \sum_{i=r+1}^m c_i^2.
\]
Số hạng thứ hai không phụ thuộc $\mat{x}$. Cực tiểu hóa số hạng đầu: $y_i = c_i/\sigma_i$
($i \leq r$), còn $y_i$ ($i > r$) tự do. Để $\norm{\mat{x}}_2$ nhỏ nhất, chọn $y_i = 0$
($i > r$). Suy ra $\hat{\mat{x}} = \mat{V}\mat{y} = \mat{V}_r\boldsymbol{\Sigma}_r^{-1}\mat{U}_r^T\mat{b}$.
\end{proof}

\begin{table}[H]
\centering
\caption{Phân loại bài toán $\mat{A}\mat{x} = \mat{b}$ và nghiệm qua SVD}
\label{tab:svd_cases}
\renewcommand{\arraystretch}{1.7}
\setlength{\tabcolsep}{5pt}
\small
\begin{tabular}{p{3.5cm}p{2.4cm}p{3.5cm}p{4.6cm}}
\toprule
\textbf{Loại hệ} & \textbf{Điều kiện} & \textbf{Nghiệm} &
\textbf{Nghiệm $\hat{\mat{x}} = \mat{A}^+\mat{b}$} \\
\midrule
Chính xác, duy nhất
  & $m = n,\; r = n$
  & $\mat{x} = \mat{A}^{-1}\mat{b}$
  & Nghiệm duy nhất \\[3pt]
Dư thừa \newline\textit{(overdetermined)}
  & $m > n,\; r = n$
  & Không có nghiệm chính xác
  & Nghiệm BPTT duy nhất \\[3pt]
Thiếu ràng buộc \newline\textit{(underdetermined)}
  & $m < n,\; r = m$
  & Vô số nghiệm
  & Nghiệm $\norm{\cdot}_2$ nhỏ nhất \\[3pt]
Tổng quát
  & $r < \min(m,n)$
  & Vô số nghiệm BPTT
  & Nghiệm BPTT có $\norm{\cdot}_2$ nhỏ nhất \\
\bottomrule
\end{tabular}
\end{table}

\subsubsection{Phân Tích Thành Phần Chính (PCA)}

Cho ma trận dữ liệu $\mat{X} \in \R^{n \times d}$ (n mẫu, d chiều), đã chuẩn hóa
zero-mean. PCA tìm $k$ hướng biến thiên lớn nhất. Cách thực hiện qua SVD:

\begin{enumerate}[noitemsep]
    \item Tính SVD: $\mat{X} = \mat{U}\boldsymbol{\Sigma}\mat{V}^T$.
    \item \textbf{Các thành phần chính} (\textit{principal components}):
          $k$ cột đầu của $\mat{V}$, tức $\mat{v}_1, \ldots, \mat{v}_k$.
    \item \textbf{Chiếu dữ liệu} xuống $k$ chiều: $\mat{Z} = \mat{X}\mat{V}_k \in \R^{n \times k}$.
    \item \textbf{Phương sai giải thích} bởi thành phần $i$:
          $\sigma_i^2 / \sum_{j=1}^r \sigma_j^2$.
\end{enumerate}

Mối liên hệ: SVD của $\mat{X}$ tương đương eigendecomposition của ma trận hiệp
phương sai $\mat{X}^T\mat{X} = \mat{V}(\boldsymbol{\Sigma}^T\boldsymbol{\Sigma})\mat{V}^T$,
nên $\mat{v}_i$ là các eigenvector của ma trận hiệp phương sai~\cite{higham2002}.

\subsubsection{Kết Luận}

Báo cáo đã trình bày đầy đủ lý thuyết phân rã SVD từ nền tảng đến ứng dụng.
Các điểm cốt lõi:

\begin{enumerate}
    \item \textbf{Tính tổng quát:} SVD tồn tại với mọi $\mat{A} \in \R^{m \times n}$
          và được xây dựng tường minh qua eigendecomposition của $\mat{A}^T\mat{A}$.
          Các giá trị kỳ dị $\sigma_i$ là duy nhất và phản ánh đầy đủ cấu trúc của $\mat{A}$.

    \item \textbf{Ý nghĩa hình học:} SVD phân tích $\mat{A}$ thành ba bước
          \textit{quay -- co giãn -- quay}; hình cầu đơn vị biến thành hyperellipsoid
          với bán trục $\sigma_i$ theo hướng $\mat{u}_i$.

    \item \textbf{Tính chất đại số:} $\rank(\mat{A})$, chuẩn spectral, chuẩn Frobenius,
          các không gian cơ bản, và giả nghịch đảo Moore--Penrose đều đọc trực tiếp
          từ SVD.

    \item \textbf{Định lý Eckart--Young:} $\mat{A}_k = \sum_{i=1}^k\sigma_i\mat{u}_i\mat{v}_i^T$
          là xấp xỉ hạng $k$ tốt nhất của $\mat{A}$ theo cả chuẩn spectral lẫn
          Frobenius, với sai số $\sigma_{k+1}$.

    \item \textbf{Ứng dụng rộng rãi:} Giải bài toán bình phương tối thiểu cho mọi
          loại hệ (overdetermined, underdetermined, rank-deficient); là nền tảng của
          PCA, nén dữ liệu, xử lý ảnh và nhiều bài toán học máy.
\end{enumerate}


